\begin{definition}
Случайный процесс $\displaystyle ( X_{t} ,\ t\geqslant 0)$ является процессом с независимыми приращениями, если $\displaystyle \forall n\in \mathbb{N} \ \forall t_{1} ,\ \dotsc ,\ t_{n}$ случайные величины $\displaystyle X_{t_{n}} -X_{t_{n-1}} ,\ \dotsc ,\ X_{t_{2}} -X_{t_{1}} ,\ X_{t_{1}}$ независимы в совокупности.
\end{definition}
\begin{note}
Случайный процесс с независимыми приращениями является непрерывным аналогом случайного блуждания.
\end{note}
\begin{theorem}[О существовании процессов с независимыми приращениями]\label{4.3} Пусть $\displaystyle \forall s,t:\ 0\leqslant s\leqslant t$ задана вероятностная мера $\displaystyle Q_{s,t}$ на $\displaystyle (\mathbb{R} ,\ \mathcal{B}(\mathbb{R}))$ с характеристической функцией $\displaystyle \phi _{s,t}$. Пусть также задана вероятностная мера $\displaystyle Q_{0}$ на $\displaystyle (\mathbb{R} ,\ \mathcal{B}(\mathbb{R}))$. Пусть также случайный процесс $\displaystyle ( X_{t} ,\ t\geqslant 0)$ является случайным процессом с независимыми приращениями и распределениями приращений
\begin{gather*}
X_{t} -X_{s}\xlongequal{d} Q_{s,t} ,\ 0\leqslant s< t,\\
X_{0}\xlongequal{d} Q_{0} .
\end{gather*}
Такой процесс существует тогда и только тогда, когда
\begin{equation}
\forall s,u,t:\ 0\leqslant s< u< t\hookrightarrow \phi _{s,t}( \tau ) =\phi _{s,u}( \tau ) \cdotp \phi _{u,s}( \tau ) .
\end{equation}
\end{theorem}
\begin{proof}
$\Longrightarrow$ $X_t - X_s = (X_t - X_u) + (X_u - X_s)$. Случайные величины в скобках независимы, поэтому $\phi_{X_s - X_t}(\tau) =\phi_{X_t - X_u}(\tau)\cdot\phi_{X_u - X_s}(\tau) = \phi_{ut}(\tau)\cdot\phi_{su}(\tau)$
\\
$\Longleftarrow$ 



Пусть $\phi_0$ - хар функция $Q_0$.Предположим, что процесс существует. Пусть также $\displaystyle 0=t_{0} < t_{1} < \dotsc < t_{n}$. Рассмотрим вектор $\displaystyle \xi =( X_{t_{n}} ,\ X_{t_{n-1}} ,\ \dotsc ,\ X_{t_{1}} ,\ X_{t_{0}})$. Найдем его характеристическую функцию. Для этого рассмотрим вектор приращений $\displaystyle \xi '=( X_{t_{n}} -X_{t_{n-1}} ,\ \dotsc ,\ X_{t_{1}} -X_{t_{0}} ,\ X_{t_{0}})$. Компоненты $\displaystyle \xi '$ независимы, следовательно, характеристическая функция $\displaystyle \xi '$ имеет вид
\begin{gather*}
\phi _{\xi '}( \lambda _{n} ,\ \dotsc ,\ \lambda _{0}) =\phi _{X_{t_{n}} -X_{t_{n-1}}}( \lambda _{n}) \cdotp \ \dotsc \ \cdotp \phi _{X_{t_{1}} -X_{0}}( \lambda _{1}) \cdotp \phi _{X_{0}}( \lambda _{0}) =\\
=\phi _{t_{n-1} ,t_{n}}( \lambda _{n}) \cdotp \ \dotsc \ \cdotp \phi _{0,t_{1}}( \lambda _{1}) \cdotp \phi _{0}( \lambda _{0}) .
\end{gather*}
Далее, заметим, что
\begin{equation*}
\xi =A\xi ',\ A=\begin{pmatrix}
1 & 1 & \dotsc  & 1\\
0 & 1 & \dotsc  & 1\\
\dotsc  & \dotsc  & \dotsc  & \dotsc \\
0 & 0 & \dotsc  & 1
\end{pmatrix} .
\end{equation*}
Тогда
\begin{gather*}
\phi _{\xi }(\overline{\lambda }) =\mathbb{E} e^{i< \overline{\lambda } ,\ \xi > } =\mathbb{E} e^{i< \overline{\lambda } ,\ A\xi '> } =\mathbb{E} e^{i\left< A^{T}\overline{\lambda } ,\ \xi '\right> } =\phi _{\xi '}\left( A^{T}\overline{\lambda }\right) =\\
=\phi _{t_{n-1} ,t_{n}}( \lambda _{n}) \cdotp \phi _{t_{n-2} ,t_{n-1}}( \lambda _{n} +\lambda _{n-1}) \cdotp \ \dotsc \ \cdotp \phi _{0}( \lambda _{n} +\ \dotsc \ +\lambda _{0}) =:\phi _{0,t_{1} ,\ \dotsc ,\ t_{n}}( \lambda _{0} ,\ \dotsc ,\ \lambda _{n}) .
\end{gather*}
Положим $\displaystyle \phi _{t_{1} ,\ \dotsc ,\ t_{n}}( \lambda _{1} ,\ \dotsc ,\ \lambda _{n}) :=\phi _{0,t_{1} ,\ \dotsc ,\ t_{n}}( 0,\ \lambda _{1} ,\ \dotsc ,\ \lambda _{n})$.

Проверим, что выполняется условие (1) из следствия для характеристических функций\\ $\displaystyle \phi _{t_{1} ,\ \dotsc ,\ t_{n}} ,\ 0\leqslant t_{1} < \dotsc < t_{n}$. Без ограничения общности проверим только ситуацию, когда есть нулевой момент времени $\displaystyle t_{0} =0$. Если $\displaystyle m=0$, то по определению $\displaystyle \phi _{t_{0} ,\ \dotsc ,\ t_{n}}( \lambda _{0} ,\ \dotsc ,\ \lambda _{n}) \big\rvert_{\lambda _{0} =0} =\phi _{t_{1} ,\ \dotsc ,\ t_{n}}( \lambda _{1} ,\ \dotsc ,\ \lambda _{n})$. Ели $\displaystyle m >0$, то по условию теоремы выполняется

\begin{gather*}
\phi _{0,\ t_{1} ,\ \dotsc ,\ t_{n}}( \lambda _{0} ,\ \lambda _{1} ,\ \dotsc ,\ \lambda _{n}) \big\rvert_{\lambda _{m} =0} =\\
=\phi _{t_{n-1} ,t_{n}}( \lambda _{n}) \cdotp \ \dotsc \ \cdotp \phi _{t_{m} ,t_{m+1}}( \lambda _{n} +\ \dotsc +\lambda _{m+1}) \cdotp \phi _{t_{m-1} ,t_{m}}( \lambda _{n} +\ \dotsc +\lambda _{m+1} +0) \cdotp \ \dotsc \ \cdotp \\
\cdotp \phi _{0,t_{1}}( \lambda _{n} +\ \dotsc \ +\lambda _{m+1} +\lambda _{m-1} +\ \dotsc \ +\lambda _{0}) =\\
=\phi _{t_{n-1} ,t_{n}}( \lambda _{n}) \cdotp \ \dotsc \ \cdotp \phi _{t_{m-1} ,t_{m+1}}( \lambda _{n} +\ \dotsc +\lambda _{m+1}) \cdotp \ \dotsc \ \cdotp \\
\cdotp \phi _{0,t_{1}}( \lambda _{n} +\ \dotsc \ +\lambda _{m+1} +\lambda _{m-1} +\ \dotsc \ +\lambda _{0}) =\\
=\phi _{0,t_{1} ,\ \dotsc ,t_{m-1} ,t_{m+1} ,\ \dotsc ,\ t_{n}}( \lambda _{0} ,\ \dotsc ,\ \lambda _{m-1} ,\ \lambda _{m+1} ,\ \dotsc ,\ \lambda _{n}) .
\end{gather*}
Таким образом, выполняется условие (1) следствия. Следовательно, по следствию существует вероятностное пространство и случайный процесс $\displaystyle ( X_{t} ,\ t\geqslant 0)$ на нем, что $\displaystyle \phi _{t_{1} ,\ \dotsc ,\ t_{n}}$ является характеристической функцией вектора $\displaystyle ( X_{t_{1}} ,\ \dotsc ,\ X_{t_{n}}) ,\ 0\leqslant t_{1} < \dotsc < t_{n}$. По построению такой процесс ялвяется процессом с независимыми приращениями, и выполняется $\displaystyle X_{t_{j}} -X_{t_{i}}\xlongequal{d} Q_{t_{j-1} ,\ t_{i}} .$
\end{proof}
\begin{definition}
     Случайный процесс $\{X_t\}_{t \geq 0}$ называется \textit{процессом Леви}, если
\begin{enumerate}
    \item $X_0 = 0$ пн
    \item $X_t$ имеет независимые приращения
    \item $\forall 0 \leq s < t, \ \forall h \geq 0$ $X_t - X_s \overset{d}{=} X_{t+h} - X_{s+h}$
\end{enumerate}
\end{definition}